\section{Combinatorial lemmas}\label{sec:comb}

In the remainder of the paper we will show that Theorem~\ref{thm:Lam-real} implies Theorem~\ref{thm:intersection-real} and more generally for affine oriented matroids, equation~(\ref{eq:Lam-determinant}) implies equation~(\ref{eq:Sq-determinant}) in the case of normal crossings.  In this section we assemble the necessary combinatorial definitions and properties of normal crossing hyperplane arrangements and affine oriented matroids.

\subsection{Normal crossing} We will make use of the following important genericity condition.
\begin{definition}\label{def:normal-crossing}
An affine oriented matroid $(\M,0)$ is \emph{normal crossing} if the
zero set $z(G)$ of every feasible covector $G$ is independent in the
underlying matroid $M$.
\end{definition}
For an affine hyperplane arrangement, this means that every nonempty
intersection of $s$ hyperplanes has codimension $s$. Equivalently, the
reduced divisor
\[
 D=\bigcup_{i\in E_0}H_i^{\CM}\subseteq\CM^r
\]
is a simple normal crossing divisor.
This condition holds for generic parameters $c_i$ in the equations $f_i(x)=\ell_i(x)+c_i$.

In oriented-matroid terminology, it corresponds to the genericity of the distinguished
element $0$, see \cite[Remark 2.3 and Section 5.2]{KowalenkoMautner}.
Normal crossing makes the local face structure Boolean, which we use
in the quadratic presentations and the Koszul complex below.

Normal crossing also implies that $M$ is the unique free coextension of $M_0$, the contraction of $M$ by $0$. This explains why
the normal crossing determinant formula involves only $M_0$.

\begin{proposition}\label{prop:normal-crossing-ranks}
If $(\M,0)$ is normal crossing, then $M$ is the free coextension of $M_0$ by $0$.
For $A\subseteq E_0$, its rank function is
\begin{equation}\label{eq:coextension-rank}
 \rk_M(A\cup\{0\})=\rk_{M_0}A+1,\qquad
 \rk_M A=\min\{\rk_{M_0}A+1,|A|\}.
\end{equation}
Its proper flats avoiding $0$ are the independent subsets of $M_0$.
Its proper flats containing $0$ are $K\cup\{0\}$ for proper flats
$K$ of $M_0$.
\end{proposition}
\begin{proof}
Every flat avoiding $0$ is the zero set of a feasible covector,
after reversing all signs if necessary. Normal crossing makes such
a flat independent. If $A$ is independent in $M_0$, it is independent
in $M$. If $A$ is dependent in $M_0$, then $0\in\cl_M A$, since
otherwise $\cl_M A$ would be an independent flat avoiding $0$.
The contraction rank formula gives \eqref{eq:coextension-rank}.
This is the rank function of the free coextension
\cite[Sections 7.2--7.3]{Oxley2011}, and the flat description follows.
\end{proof}

To compute the weights $w_F$ when $(\M,0)$ is normal crossing, we
use the following identity for the beta invariant.
For a matroid $N$ on $E_0$, write $\widehat N$ for its free coextension
by $0$. Its rank function is given by \eqref{eq:coextension-rank}, with
$(M,M_0)$ replaced by $(\widehat N,N)$.

\begin{lemma}\label{lem:coextension-beta}
For any matroid $N$,
\[
 \beta(\widehat N)=\mup(N^*).
\]
\end{lemma}
\begin{proof}
This follows from Brylawski's free extension formula for the beta
invariant and duality
\cite[Remarks 6.15(3)(b) and (d), p.~241]{Brylawski1982}.
\end{proof}

\begin{lemma}\label{lem:free-coextension}
Let $(\M,0)$ be normal crossing and $M_0=M/0$. For a nonempty proper
flat $F$ of $M$, 
\[
 w_F= \beta(M^F)\beta(M_F) = 
 \begin{cases}
  \beta(M_{\{i\}}),&F=\{i\},\quad i\in E_0,\\
  0,&0\notin F,\quad |F|>1,\\
  \beta(M_{0,K})\mup((M_0^K)^*),&F=K\cup\{0\}.
 \end{cases}
\]
\end{lemma}
\begin{proof}
By Proposition~\ref{prop:normal-crossing-ranks}, for a flat avoiding
$0$, $M^F$ is free and its beta invariant vanishes unless
$F$ is a singleton. For $F=K\cup\{0\}$,
\[
 M^F=\widehat{M_0^K},\qquad M_F=M_{0,K}.
\]
Apply Lemma~\ref{lem:coextension-beta} to the restriction.
\end{proof}

\subsection{Face lattices} The following are standard observations for affine hyperplane
arrangements. We use the covector axioms and the cited results to
obtain their oriented-matroid versions.

We write
$\bar P\cap\bar Q\cap H_J\ne\varnothing$ if there is a
feasible $G\leq P,Q$ with $J\subseteq z(G)$. Moreover, we define
\[
 \Delta_{PQ}=\{J\subseteq E_0\mid\bar P\cap\bar Q\cap H_J\ne\varnothing\}.
\]
Thus $\Delta_{PQ}$ records the collections of hyperplanes whose
intersection meets the common face of $P$ and $Q$.
It is closed under taking subsets.


Near a normal crossing face, either side of each hyperplane containing
the face can be chosen independently. This gives a cube of feasible
regions. Conversely, if all sign choices on $S$ occur while the signs
outside $S$ are fixed, convexity gives a common face on the hyperplanes
in $S$ with those fixed signs.

\begin{lemma}\label{lem:feasible-cubes}
Let $P,Q\in\cF$, $J\subseteq E_0$, and $S=\Sep(P,Q)\cup J$.
If $P^K\in\cF$ for every $K\subseteq S$, then setting the entries of
$P$ on $S$ to zero gives a feasible covector. In particular,
$J\in\Delta_{PQ}$.
If $(\M,0)$ is normal crossing, $J\in\Delta_{PQ}$ implies that every
$P^K$, $K\subseteq S$, is feasible.
\end{lemma}
\begin{proof}
The first assertion follows by strong elimination and induction on $|S|$.

Conversely, $J\in\Delta_{PQ}$ gives a feasible $G\leq P,Q$ with
$J\subseteq z(G)$, and hence $S\subseteq z(G)$. If $(\M,0)$ is normal
crossing, the restriction to $z(G)$ is free, so every sign assignment there
extends to a covector $Y$. The covector $G\circ Y$ fills those zero
entries and preserves all other signs, including the sign at $0$.
Choosing the signs on $z(G)\setminus S$ to agree with $P$ gives every
vertex $P^K$ of the cube.
\end{proof}

\begin{corollary}\label{cor:face-subsets}
Assume that $(\M,0)$ is normal crossing. The zero-set map is an order-reversing
bijection between the feasible faces $G\leq P$ and the simplices of
$\Delta_{PP}$. The face indexed by $J$ has zero entries on $J$ and
the signs of $P$ elsewhere. If $P,Q$ have a nonempty common face,
the covector of $\bar P\cap\bar Q$ has zero set $\Sep(P,Q)$.
\end{corollary}
\begin{proof}
Apply Lemma~\ref{lem:feasible-cubes} with $Q=P$ for the first
assertion. A face below $P$ is uniquely determined by its zero set.
For the second assertion, apply the lemma with $J=\varnothing$.
The resulting covector is the largest common face of $P,Q$.
\end{proof}

Assume that $(\M,0)$ is normal crossing. For $P,R\in\cF$, put
\[
 \begin{aligned}
 L_P^R&=\{z(G)\mid G\leq P\text{ feasible},\ z(G)\subseteq\Sep(P,R)\}\\
      &=\{J\in\Delta_{PP}\mid J\subseteq\Sep(P,R)\}.
 \end{aligned}
\]
The equality follows from Corollary~\ref{cor:face-subsets}, so this
is a simplicial complex.
Its vertices label the facets of $\bar P$ separating $P$ from $R$.
For hyperplanes these are the facets visible from a point of $R$.
Their union projects onto a convex set, and their nonempty
intersections are convex. The nerve lemma therefore gives
contractibility when $P\ne R$.

\begin{lemma}\label{lem:separating-facets}
If $(\M,0)$ is normal crossing, then $L_P^R$ is contractible for $P\ne R$.
\end{lemma}
\begin{proof}
The face-filter lemma \cite[Lemma 4.3.11]{OrientedMatroids1999}
says that for topes $P\ne\pm R$ the poset
\[
 \{G\mid 0<G<P,\ z(G)\subseteq\Sep(P,R)\}
\]
has contractible order complex. Here the order complex has a vertex
for each covector and a simplex for each chain of covectors.
Since $P(0)=R(0)=+$, these covectors are feasible and $P\ne\pm R$.
The zero-set map identifies this poset with the nonempty simplices of
$L_P^R$, ordered by reverse inclusion. Its order complex is therefore
the barycentric subdivision of $L_P^R$, which has the same underlying
space. The assertion follows.
\end{proof}

The last lemma converts a weighted count of hyperplane intersections
meeting a face into a count of its bounded faces. For a ray in
dimension one, it reads $(1-t)+t=1$, with only the endpoint contributing
on the right.

\begin{lemma}\label{lem:face-enumeration}
If $(\mathcal M,0)$ is normal crossing, for $P,Q\in\cF$ and $F=\bar P\cap\bar Q$,
\[
 \sum_{J\in\Delta_{PQ}}t^{|J|}(1-t)^{r-|J|}=h(F_\bd,t).
\]
\end{lemma}
\begin{proof}
We may assume $F\ne\varnothing$.
Expanding $t^{|J|}=(1+(t-1))^{|J|}$ gives
\[
 \sum_{J\in\Delta_{PQ}}t^{|J|}(1-t)^{r-|J|}
 =\sum_{I\in\Delta_{PQ}}c_I(t-1)^{r-|I|},
\text{ where }
 c_I=\sum_{\substack{J\in\Delta_{PQ}\\I\subseteq J}}(-1)^{r-|J|}.
\]

If $z(F)\not\subseteq I$, choose $i\in z(F)\setminus I$. Since
$F\subseteq H_i$, we have $
 F\cap H_J=F\cap H_{J\cup\{i\}}.$
Hence the terms in $c_I$ cancel in pairs indexed by $J$
and $J\cup\{i\}$, with $i\notin J$, giving $c_I=0$.

If $z(F)\subseteq I$, Corollary~\ref{cor:face-subsets} gives a unique
feasible face $G\leq F$ with $z(G)=I$. The sets $J\in\Delta_{PQ}$
containing $I$ correspond to the feasible faces $H\leq G$.
For these faces, normal crossing gives $\dim H=r-|z(H)|$, so
\[
 c_I=\sum_{\substack{H\leq G\\H \text{ feasible}}}(-1)^{\dim H}.
\]
Denote by $G^\infty\leq G$ the largest covector below $G$ with
$G^\infty(0)=0$, obtained by composing all such covectors. Then
$H\leq G$ is not feasible if and only if $H\leq G^\infty$.
Moreover, $G$ is bounded if and only if $G^\infty=\mathbf0$.

We will make use of the Euler relation for covectors, see 
\cite[Corollary 4.3.8]{OrientedMatroids1999}, which states that for every nonzero covector $G'$
\[
 \sum_{\mathbf0<H\leq G'}(-1)^{\dim H}=1.
\]
Applying this to $G'=G$ and splitting the sum, we obtain
\[
 \begin{aligned}
  1 &=\sum_{\mathbf0<H\leq G}(-1)^{\dim H}  = c_I
    +\sum_{\mathbf0<H\leq G^\infty }
       (-1)^{\dim H}.
 \end{aligned}
\]
If $G$ is bounded $G^\infty=\mathbf{0}$ and the second summand is zero, so $c_I=1$. Otherwise, the Euler identity for $G'=G^\infty$ applies to the second summand so that $c_I=0$.

Consequently, we obtain
\[
 \sum_{J\in\Delta_{PQ}}t^{|J|}(1-t)^{r-|J|}
 =\sum_{\substack{G\leq F\\G\text{ bounded}}}(t-1)^{\dim G}
 =h(F_\bd,t)\qedhere
\]
\end{proof}
